\documentclass[12pt]{article}
\usepackage[utf8]{inputenc}
\usepackage[T1]{fontenc}
\usepackage{geometry}
\geometry{a4paper, margin=2.5cm}
\usepackage{setspace}
\onehalfspacing
\usepackage{hyperref}
\usepackage{listings}
\usepackage{xcolor}
\usepackage{longtable}
\usepackage{tabularx}
\usepackage{booktabs}
\usepackage{amsmath}
\usepackage{float}
\usepackage{array}
\usepackage{ragged2e}

\lstset{
  basicstyle=\ttfamily\footnotesize,
  breaklines=true,
  breakatwhitespace=false,
  breakautoindent=true,
  postbreak=\mbox{\textcolor{red}{$\hookrightarrow$}\space},
  backgroundcolor=\color{gray!10},
  frame=single,
  numbers=left,
  numberstyle=\tiny,
  captionpos=b,
  xleftmargin=0.5em,
  xrightmargin=0.5em
}

\newcolumntype{L}{>{\RaggedRight\arraybackslash}X}
\newcolumntype{P}[1]{>{\RaggedRight\arraybackslash}p{#1}}

\newcommand{\tabcol}[1]{>{\RaggedRight\arraybackslash}p{#1}}
\newlength{\descrwidth}
\setlength{\descrwidth}{3.8cm}

\title{Documento de Procedencia de Datos\\\small{Sistema de Predicci\'on Temprana de HABs}}
\author{}
\date{}

\begin{document}
\maketitle
\thispagestyle{empty}
\newpage

\section{Fuentes de Datos}

\subsection{Sentinel-2 (ESA Copernicus)}
\begin{tabular}{P{\descrwidth} P{\dimexpr\textwidth-\descrwidth-4\tabcolsep\relax}}
\toprule
\textbf{Agente:} & ESA Copernicus Data Space \\
\midrule
\textbf{Producto:} & S2MSI2A (Bottom-of-Atmosphere, BOA) \\
\midrule
\textbf{Bandas:} & B2 (493~nm), B3 (560~nm), B4 (665~nm), B5 (705~nm), B8 (842~nm) \\
\midrule
\textbf{Resoluci\'on espacial:} & 10--20~m (remuestreadas a 20~m) \\
\midrule
\textbf{Resoluci\'on temporal:} & 5 d\'ias \\
\midrule
\textbf{Proyecci\'on:} & UTM zona 17N (EPSG:32617) \\
\midrule
\textbf{Formato:} & GeoTIFF, 16-bit sin signo \\
\midrule
\textbf{Periodo:} & 2023 a 2026 \\
\midrule
\textbf{Regiones:} & Okeechobee, TampaBay, Golfo de Fonseca, Lago de Yojoa, Embalse El Caj\'on \\
\midrule
\textbf{Directorio base:} & \texttt{imagenes/\{Region\}/\{A\~no\}/} \\
\bottomrule
\end{tabular}

\textbf{Preprocesamiento:} Las bandas se normalizan dividiendo por 10000 cuando el valor m\'aximo supera 1.5. A partir de las 5 bandas se calculan 6 \'indices espectrales: NDCI, NDVI, FAI, B5/B4, B3/B2, B8/B3.

\subsection{ERA5-Land (Copernicus Climate Data Store)}
\begin{tabular}{P{\descrwidth} P{\dimexpr\textwidth-\descrwidth-4\tabcolsep\relax}}
\toprule
\textbf{Agente:} & ECMWF / Copernicus CDS \\
\midrule
\textbf{API:} & CDS API v3 (\texttt{cdsapi} Python) \\
\midrule
\textbf{Dataset:} & \texttt{reanalysis-era5-land} \\
\midrule
\textbf{Resoluci\'on espacial:} & $0.1^\circ \times 0.1^\circ$ ($\sim$9~km) \\
\midrule
\textbf{Resoluci\'on temporal:} & 4 horaria (00, 06, 12, 18 UTC) \\
\midrule
\textbf{Extensi\'on:} & Lat 31.25$^\circ$N a 14.5$^\circ$N, Lon 88.5$^\circ$W a 79.75$^\circ$W \\
\midrule
\textbf{Variables instant\'aneas:} & t2m, u10, v10, sp \\
\midrule
\textbf{Variables acumuladas:} & ssrd, tp \\
\midrule
\textbf{Periodo:} & 2023-01 a 2026-05 \\
\midrule
\textbf{Archivos:} & 211 NetCDF, 219~MB total \\
\midrule
\textbf{Directorio:} & \texttt{era5\_temp\_nc/} \\
\bottomrule
\end{tabular}

\textbf{Extracci\'on:} Se extrae el valor a las 12:00 UTC (\'indice 2 del bloque 4-horario) para cada muestra. Para el an\'alisis de predicci\'on temprana se extraen adem\'as valores en t-1, t-3, t-5 y t-7 d\'ias desde la misma cuadr\'icula ERA5.

\subsection{MODIS Aqua CHL (NASA Ocean Color)}
\begin{tabular}{P{\descrwidth} P{\dimexpr\textwidth-\descrwidth-4\tabcolsep\relax}}
\toprule
\textbf{Agente:} & NASA Ocean Color Web (\texttt{oceancolor.gsfc.nasa.gov}) \\
\midrule
\textbf{Producto:} & MODIS Aqua L3m CHL mensual, 4~km \\
\midrule
\textbf{Algoritmo:} & OC3M/OCI combinado (chlor\_a) \\
\midrule
\textbf{Periodo:} & 2023-01 a 2026-05 \\
\midrule
\textbf{Regiones:} & Golfo de Fonseca, Lago de Yojoa, Embalse El Caj\'on \\
\midrule
\textbf{Formato:} & NetCDF \\
\midrule
\textbf{Cobertura v\'alida:} & Golfo de Fonseca: 35/41 meses con CHL $>$ 0 \\
& Lago de Yojoa: 6/41 meses con CHL $>$ 0 \\
& El Caj\'on: 0/41 meses (embalse angosto, no resuelto por 4~km) \\
\midrule
\textbf{Nota:} & 2026 no produjo valores v\'alidos en ninguna estaci\'on hondure\~na \\
\bottomrule
\end{tabular}

\subsection{WQP Chlorophyll (EPA/USGS)}
\begin{tabular}{P{\descrwidth} P{\dimexpr\textwidth-\descrwidth-4\tabcolsep\relax}}
\toprule
\textbf{Agente:} & Water Quality Portal (\texttt{www.waterqualitydata.us}) \\
\midrule
\textbf{API:} & REST (\texttt{dataRetrieval} R / \texttt{wqptoolbox}) \\
\midrule
\textbf{Par\'ametro:} & Chlorophyll a (00665, 70960) \\
\midrule
\textbf{Unidades:} & ug/L y mg/m$^3$ \\
\midrule
\textbf{Organizaciones:} & STF\_WQX (South Florida Water Management), UWFCEDB (Univ. West Florida) \\
\midrule
\textbf{Periodo:} & 2023 a 2026 \\
\midrule
\textbf{Registro v\'alido:} & 59 muestras en 2025, 7 en 2026 (Panhandle de Florida) \\
\midrule
\textbf{Archivo:} & \texttt{datasets/florida/\-florida\_chlorophyll\-\_wqp\_2025\_2026.csv} \\
\midrule
\textbf{Nota:} & Las muestras de 2026 (Panhandle) no se emparejaron con Sentinel-2 por falta de im\'agenes en esa zona \\
\bottomrule
\end{tabular}

\section{Scripts y Pipeline}

Todos los scripts se ejecutaron en Windows 11 con PowerShell 5.1 y Python 3.13, ocupando GPU NVIDIA CUDA para el entrenamiento de redes neuronales.

\subsection{\texttt{calibrar\_chl\_s2.py}}
\textbf{Prop\'osito:} Entrenar modelo RandomForest para estimar clorofila desde Sentinel-2. Usa 193 muestras reales con clorofila in-situ.

\textbf{C\'odigo:}
\begin{lstlisting}[language=Python]
F_CHL = ["B2","B3","B4","B5","B8","NDCI","NDVI","FAI",
         "B5_B4_ratio","B3_B2_ratio","B8_B3_ratio"]
chl_log = np.log10(np.maximum(chl, 0.1))

rf_chl = RandomForestRegressor(
    n_estimators=300, max_depth=10,
    min_samples_leaf=3, random_state=RND, n_jobs=-1
)
rf_chl.fit(X, chl_log)

pred_chl = 10 ** rf_chl.predict(X)
pickle.dump({"model": rf_chl, "features": F_CHL}, f)
\end{lstlisting}

\textbf{Ejecuci\'on:}
\begin{lstlisting}
python scripts/calibrar_chl_s2.py
\end{lstlisting}

\textbf{M\'etricas obtenidas:}
\begin{itemize}
  \item Pearson r: 0.85 (p $<$ 0.001)
  \item Spearman $\rho$: 0.91 (p $<$ 0.001)
  \item MAE: 6.70~ug/L
\end{itemize}
\textbf{Output:} \texttt{modelo\_final\_2023/rf\_chl\_estimator.pkl}

\subsection{\texttt{apply\_chl\_2026.py}}
\textbf{Prop\'osito:} Aplicar el estimador de clorofila a im\'agenes Sentinel-2 de 2026 y extraer ERA5. Genera 1,342 muestras sint\'eticas.

\textbf{C\'odigo (carga ERA5 y predicci\'on):}
\begin{lstlisting}[language=Python]
def load_era5_month(y, m):
    vt = ds_i.valid_time.values
    mask = np.array([t.hour == 12
                     for t in pd.DatetimeIndex(vt)])
    t2m = ds_i.t2m.values[mask]
    # ... same for sp, u10, v10, ssrd, tp
    wspd = np.sqrt(u10**2 + v10**2)
    return {"t2m": t2m, "sp": sp, "wspd": wspd, ...}

def get_era5_bulk(fecha, lat, lon, cache):
    day = fecha.day - 1
    ilat = np.argmin(np.abs(d["lats"] - lat))
    ilon = np.argmin(np.abs(d["lons"] - lon))
    return {"temp_air_2m": float(d["t2m"][day,ilat,ilon]),
            "solar_radiation": float(d["ssrd"][day,ilat,ilon]),
            "precipitation": float(d["tp"][day,ilat,ilon]),
            "wind_speed_10m": float(d["wspd"][day,ilat,ilon]),
            ...}
\end{lstlisting}

\textbf{C\'odigo (submuestreo de p\'ixeles):}
\begin{lstlisting}[language=Python]
step = max(1, int(np.sqrt(H * W / 5000)))
b_sub = b[:, ::step, ::step]
bf = b_sub.reshape(5, -1).T
water = bf.sum(1) > 0
idxs = np.random.choice(water.sum(), SAMPLES_PER_TILE, replace=False)

for idx in idxs:
    utm_x, utm_y = tr * (col + step//2, row + step//2)
    # UTM -> WGS84
    lon, lat = rio_transform(
        crs_src, {"init": "EPSG:4326"}, [utm_x], [utm_y])
    chl_est = float(10 ** rf_chl.predict(
        feats.reshape(1, -1))[0])
    era = get_era5_bulk(fecha, lat, lon, era_cache)
\end{lstlisting}

\textbf{Ejecuci\'on:}
\begin{lstlisting}
python scripts/apply_chl_2026.py
\end{lstlisting}
\textbf{Output:} \texttt{datasets/samples\_2026\_s2\_chl.csv} (1,342 muestras)

\subsection{\texttt{lead\_time\_era5.py}}
\textbf{Prop\'osito:} Extraer ERA5 de t-1, t-3, t-5, t-7 d\'ias para cada muestra, demostrando predicci\'on temprana.

\textbf{C\'odigo:}
\begin{lstlisting}[language=Python]
LEADS = [("t0", 0), ("t-1", -1), ("t-3", -3),
         ("t-5", -5), ("t-7", -7)]
ERA_VARS = ["temp_air_2m","solar_radiation",
            "precipitation","wind_speed_10m",
            "wind_u_10m","wind_v_10m",
            "surface_pressure"]

for lead_name, lead_days in LEADS:
    lead_fecha = fecha + pd.Timedelta(days=lead_days)
    era = get_era5_at(lead_fecha, lat, lon)
    for v in ERA_VARS:
        entry[f"{v}_{lead_name}"] = era[v]
\end{lstlisting}

\textbf{Ejecuci\'on:}
\begin{lstlisting}
python scripts/lead_time_era5.py
\end{lstlisting}
\textbf{Output:} \texttt{datasets/lead\_time\_era5.csv} (ERA5 multi-lead para 1,342 muestras)

\subsection{\texttt{clean\_concat.py}}
\textbf{Prop\'osito:} Unir datasets hist\'oricos (2023--2025) con las nuevas muestras de 2026 en un \'unico CSV.

\textbf{C\'odigo:}
\begin{lstlisting}[language=Python]
original = full_ds[(full_ds["fecha"].dt.year >= 2023) &
                   (full_ds["fecha"].dt.year <= 2025) &
                   full_ds["temp_air_2m"].notna()].copy()
s26 = pd.read_csv("datasets/samples_2026_s2_chl.csv")
s26 = s26[s26["temp_air_2m"].notna()].copy()

full = pd.concat([original_out, s26_out], ignore_index=True)
full = full.sort_values("fecha").reset_index(drop=True)
full.to_csv("datasets/dataset_entrenamiento_completo_2023_2026.csv", index=False)
\end{lstlisting}

\textbf{Ejecuci\'on:}
\begin{lstlisting}
python scripts/clean_concat.py
\end{lstlisting}
\textbf{Output:} \texttt{datasets/dataset\_entrenamiento\_completo\_2023\_2026.csv}

\section{C\'odigo de Modelos (Notebook)}
El notebook \texttt{Modelo\_tesis.ipynb} contiene todo el pipeline de entrenamiento, evaluaci\'on y generaci\'on de mapas.

\subsection{HABNet1Dv2 (Arquitectura)}
\begin{lstlisting}[language=Python]
class HABNet1Dv2(torch.nn.Module):
    def __init__(self, nf, h=128, nr=2, d=0.2):
        super().__init__()
        self.input_proj = nn.Sequential(
            nn.Linear(nf,h), nn.BatchNorm1d(h),
            nn.GELU(), nn.Dropout(d))
        self.res_blocks = nn.ModuleList([
            nn.Sequential(nn.Linear(h,h), nn.BatchNorm1d(h),
                nn.GELU(), nn.Dropout(d),
                nn.Linear(h,h), nn.BatchNorm1d(h))
            for _ in range(nr)])
        self.se = nn.Sequential(
            nn.Linear(h,max(h//4,8)), nn.ReLU(),
            nn.Linear(max(h//4,8),h), nn.Sigmoid())
        self.classifier = nn.Sequential(
            nn.Linear(h,h//2), nn.GELU(),
            nn.Dropout(d*.5), nn.Linear(h//2,1))
    def forward(self, x):
        x = self.input_proj(x)
        for b in self.res_blocks:
            x = x + b(x)  # residual connection
        return self.classifier(x * self.se(x))
\end{lstlisting}

\subsection{Entrenamiento (5-Fold CV + Hold-Out)}
\begin{lstlisting}[language=Python]
skf = StratifiedKFold(5, shuffle=True, random_state=RND)
for f, (ti, vi) in enumerate(skf.split(X, y)):
    dl_tr = DataLoader(TensorDataset(
        tensor(X_tr_s), tensor(y_tr).unsqueeze(1)), 32, True)
    dl_te = DataLoader(TensorDataset(
        tensor(X_te_s), tensor(y_te).unsqueeze(1)), 32, False)
    m = HABNet1Dv2(len(F)).to(DEVICE)
    o = optim.AdamW(m.parameters(), lr=1e-3, weight_decay=1e-4)
    s = CosineAnnealingLR(o, 100, 1e-5)
    c = FocalLoss(a=.25, g=2, pw=tensor([pw]).to(DEVICE))
    for ep in range(200):
        for xb, yb in dl_tr:
            o.zero_grad()
            c(m(xb.to(DEVICE)), yb.to(DEVICE)).backward()
            clip_grad_norm_(m.parameters(), 1)
            o.step()
        # early stopping si no mejora en 25 epochs
        if va > ba: ba = va; es = 0
        else: es += 1
        if es >= 25: break
\end{lstlisting}

\subsection{XGBoost (Configuraci\'on)}
\begin{lstlisting}[language=Python]
pw_xgb = int((y_tr == 0).sum()) / max(int(y_tr.sum()), 1)
xgb_m = xgb.XGBClassifier(
    n_estimators=300, max_depth=6, learning_rate=.05,
    subsample=.8, colsample_bytree=.8,
    scale_pos_weight=pw_xgb, random_state=RND,
    device='cuda', tree_method='hist', verbosity=0)
xgb_m.fit(X_tr_s, y_tr)
xgb_m.save_model("modelo_final_2023/xgb_model.json")
\end{lstlisting}

\subsection{Generaci\'on de Mapas (\texttt{proc\_tif})}
\begin{lstlisting}[language=Python]
def proc_tif(path, dom, st=20):
    with rasterio.open(path) as s:
        b = s.read().astype('float32') / 10000.
        pr, tr = s.profile, s.transform
    if st > 1: b = b[:, ::st, ::st]
    bf = b.reshape(5, -1).T
    v = bf.sum(1) > 0  # solo pixeles con agua
    Xs = feats({sb: bv[:,i] for i,sb in enumerate(S2)})
    if dom != 'Florida': Xs += shift  # domain adaptation
    Xf = np.column_stack([Xs, np.full((len(Xs),len(ERA_M)),ERA_M)])
    p = xgb_m.predict_proba(sc.transform(Xf.astype('float64')))[:,1]
    pm = np.zeros((H, W), 'float32')
    pm.flat[v] = p
    return pm, pr, tr
\end{lstlisting}

\textbf{Domain Adaptation:} El archivo \texttt{yojoa\_s2\_matched.csv} contiene firmas espectrales de muestras en Lago de Yojoa. Se calcula \texttt{spectral\_shift.npy} como la diferencia entre la media espectral de Florida y la media de Yojoa, y se suma a las bandas de Honduras durante la generaci\'on de mapas. Adicionalmente existe \texttt{spectral\_adaptation.pkl} (z-score), generado en experimentos previos. El m\'etodo usado en el pipeline final es el desplazamiento aditivo (\texttt{spectral\_shift.npy}).

\section{Validaciones}

\subsection{Validaci\'on Cruzada (5-Fold Stratified)}
\begin{table}[H]
\centering
\begin{tabular}{cr}
\toprule
Fold & AUC \\
\midrule
1 & 0.9813 \\
2 & 0.9844 \\
3 & 0.9800 \\
4 & 0.9795 \\
5 & 0.9811 \\
\midrule
\text{Media} & 0.9813 $\pm$ 0.002 \\
\bottomrule
\end{tabular}
\end{table}

\subsection{Hold-Out (80/20)}
\begin{table}[H]
\centering
\begin{tabular}{lccc}
\toprule
Modelo & AUC & F1 & Accuracy \\
\midrule
HABNet1Dv2 & 0.9861 & 0.8871 & 0.9419 \\
XGBoost & \textbf{0.9948} & \textbf{0.9498} & \textbf{0.9730} \\
\bottomrule
\end{tabular}
\end{table}

\subsection{Correlaci\'on con Clorofila in-situ (n = 2,424)}
\begin{align*}
\text{HABNet1Dv2:} &\quad r = 0.7314,\ \rho = 0.6247 \\
\text{XGBoost:} &\quad r = 0.7471,\ \rho = 0.7709 \\
\text{CHL HAB:} &\quad 22.8\ \mu\text{g/L} \\
\text{CHL No-HAB:} &\quad 5.4\ \mu\text{g/L}
\end{align*}

\subsection{Estabilidad Temporal (AUC por a\~no)}
\begin{table}[H]
\centering
\begin{tabular}{crrl}
\toprule
A\~no & Muestras & HAB & XGBoost AUC \\
\midrule
2023 & 118 & 51 & 0.9918 \\
2024 & 53 & 2 & 1.0000 \\
2025 & 1194 & 306 & 1.0000 \\
2026 & 1044 & 281 & 0.9985 \\
\bottomrule
\end{tabular}
\end{table}

\subsection{Predicci\'on Temprana (Lead-Time)}
Evaluaci\'on del modelo XGBoost entrenado, substituyendo ERA5 real por ERA5 de d\'ias anteriores. Las bandas espectrales S2 se mantienen fijas.

\begin{table}[H]
\centering
\begin{tabular}{crrrr}
\toprule
Lead & n & AUC HABNet & AUC XGBoost & Pearson XGB \\
\midrule
t0 (real) & 1342 & 0.9385 & 0.9747 & 0.7960 \\
t-1 & 1342 & 0.9370 & 0.9737 & 0.7896 \\
t-3 & 1342 & 0.9377 & 0.9732 & 0.7976 \\
t-5 & 1342 & 0.9351 & 0.9747 & 0.7870 \\
t-7 & 1342 & 0.9401 & \textbf{0.9754} & 0.7974 \\
\bottomrule
\end{tabular}
\caption{El AUC de XGBoost no se degrada al usar ERA5 de hasta 7 d\'{\i}as antes. Esto valida que el sistema puede operar con pron\'ostico clim\'atico en lugar de ERA5 real.}
\end{table}

\subsection{Perturbaci\'on ERA5: real vs promedio}
\begin{align*}
\text{Real:} &\quad \text{HABNet AUC}=0.9956,\ \text{XGB AUC}=0.9993 \\
\text{Promedio:} &\quad \text{HABNet AUC}=0.4824,\ \text{XGB AUC}=0.9922 \\
\text{Cambio medio HABNet:} &\quad 0.3893 \\
\text{Cambio medio XGBoost:} &\quad 0.0447
\end{align*}
El XGBoost es robusto al uso de ERA5 promedio, mientras que HABNet1Dv2 colapsa. Por esta raz\'on, XGBoost es el modelo recomendado para la generaci\'on de mapas.

\section{Dataset Final}

\begin{table}[H]
\centering
\sloppy
\begin{tabular}{P{\descrwidth} P{\dimexpr\textwidth-\descrwidth-4\tabcolsep\relax}}
\toprule
Muestras totales & 2,424 (645 HAB, 26.6\%) \\
\midrule
Eventos HAB & 645 (26.6\%) \\
\midrule
Periodo & 2023-01 a 2026-05 \\
\midrule
Variables espectrales & 11 (5 bandas + 6 \'indices) \\
\midrule
Variables clim\'aticas & 7 (ERA5) \\
\midrule
Origenes & golfo\_fonseca (1,422), florida (726), lago\_de\_yojoa (140), caj\'on (136) \\
\midrule
Archivo & \texttt{datasets/dataset\_entrenamiento\_completo\_2023\_2026.csv} \\
\bottomrule
\end{tabular}
\end{table}

\section{Estructura de Archivos}
{\scriptsize
\begin{lstlisting}
C:\Users\JC\Desktop\Tesis\
+-- Modelo_tesis.ipynb
+-- datasets/
|   +-- dataset_entrenamiento_completo_2023_2026.csv
|   +-- lead_time_era5.csv
|   +-- samples_2026_s2_chl.csv
|   +-- honduras_modis_chl_monthly.csv
|   +-- florida/
|       +-- ..._wqp_2025_2026.csv
|       +-- florida_s2_matched.csv
+-- imagenes/
|   +-- {Okeechobee,TampaBay,Cajon,
|        Golfo_Fonseca,Lago de Yojoa}/
|       +-- {2023,2024,2025,2026}/*.tif
+-- era5_temp_nc/   (211 archivos)
+-- modelo_final_2023/
|   +-- best_model.pth           (357 KB) HABNet final
|   +-- habnet1d_v2_temporal.pt  (360 KB) HABNet temp
|   +-- habnet1d_v2_final.pt     (357 KB) HABNet final alt
|   +-- xgb_model.json           (815 KB) XGBoost final
|   +-- xgboost_final.pkl        (278 KB) XGBoost final alt
|   +-- xgb_temporal.pkl         (443 KB) XGBoost temporal
|   +-- scaler.pkl               StandardScaler
|   +-- scaler_final.pkl         StandardScaler alt
|   +-- scaler_temporal.pkl      StandardScaler temp
|   +-- rf_chl_estimator.pkl     (1.5 MB) RF CHL
|   +-- spectral_shift.npy       Domain adaptation shift
|   +-- spectral_adaptation.pkl  Domain adaptation z-score
|   +-- summary_final.json       Metrics final
|   +-- summary_temporal.json    Metrics temporal
|   +-- roc_final.png, roc_temporal.png
+-- mapas_finales/
|   +-- *_hab_final.tif   (867 mapas)
+-- scripts/
    +-- calibrar_chl_s2.py   # RF CHL estimator
    +-- apply_chl_2026.py    # 2026 sample generation
    +-- lead_time_era5.py    # Multi-lead ERA5
    +-- clean_concat.py      # Dataset consolidation
    +-- concat_2026.py       # Alt. concatenation script
\end{lstlisting}
}

\section*{Resumen de Ejecuci\'on}

Todos los scripts se ejecutaron en el siguiente orden:
\begin{enumerate}
  \item \texttt{calibrar\_chl\_s2.py} (RF CHL, Pearson r=0.85)
  \item \texttt{apply\_chl\_2026.py} (1,342 muestras 2026)
  \item \texttt{lead\_time\_era5.py} (ERA5 t-1 a t-7)
  \item \texttt{clean\_concat.py} (dataset final 2,424 muestras)
  \item \texttt{Modelo\_tesis.ipynb} (entrenamiento, mapas, predicci\'on temprana)
\end{enumerate}

\end{document}
